\section*{Bab \ref{ch:g10:algebra} --- Aljabar: Persamaan dan Pertidaksamaan}

\begin{solution}{exo:g10:algebra:1}
$3(2x-5) - 2(x-7) = 6x - 15 - 2x + 14 = 4x - 1$.

$(x+4)(2x-1) = 2x^2 - x + 8x - 4 = 2x^2 + 7x - 4$.

$(3x-2)^2 = 9x^2 - 12x + 4$.

$(5-x)(5+x) = 25 - x^2$.
\end{solution}

\begin{solution}{exo:g10:algebra:2}
$x^2 - 49 = (x+7)(x-7)$ (selisih dua kuadrat).

$x^2 + 10x + 25 = (x+5)^2$.

$7x^2 - 14x = 7x(x - 2)$ (faktor persekutuan).

$16x^2 - 8x + 1 = (4x - 1)^2$.
\end{solution}

\begin{solution}{exo:g10:algebra:3}
$4x + 9 = x - 3$: $3x = -12$, jadi $x = -4$.

$\frac{2x-1}{3} = 5$: $2x - 1 = 15$, jadi $x = 8$.

$2(x-3) = 2x + 1$ terjabar menjadi $2x - 6 = 2x + 1$, yaitu $-6 = 1$: salah
untuk setiap $x$. Tidak ada penyelesaian --- kedua ruasnya menyatakan garis
sejajar yang tidak pernah bertemu.
\end{solution}

\begin{solution}{exo:g10:algebra:4}
$(x+5)(2x-8) = 0$: $x = -5$ atau $x = 4$.

$x^2 - 16 = (x-4)(x+4) = 0$: $x = 4$ atau $x = -4$.

$x^2 + 6x + 9 = (x+3)^2 = 0$: penyelesaian tunggalnya $x = -3$.
\end{solution}

\begin{solution}{exo:g10:algebra:5}
$3x - 5 \leq 7$: $3x \leq 12$, jadi $x \leq 4$: himpunan penyelesaiannya
$\intoc{-\infty}{4}$.

$-4x + 1 > 9$: $-4x > 8$, membagi dengan $-4$ membalik arah pertidaksamaan:
$x < -2$: himpunan penyelesaiannya $\intoo{-\infty}{-2}$.

$2(x-1) \geq 5x + 4$: $2x - 2 \geq 5x + 4$, jadi $-3x \geq 6$, jadi
$x \leq -2$: himpunan penyelesaiannya $\intoc{-\infty}{-2}$.
\end{solution}

\begin{solution}{exo:g10:algebra:6}
Faktor persekutuannya adalah $(x-1)$:
\[
E = (x-1)\bigl[(x+3) - (2x - 5)\bigr] = (x-1)(x + 3 - 2x + 5)
= (x-1)(8 - x).
\]
Lalu $E = 0$ ketika $x = 1$ atau $x = 8$ (aturan hasil kali nol).
\end{solution}

\begin{solution}{exo:g10:algebra:7}
$x^2 = 5x$: $x^2 - 5x = x(x - 5) = 0$, jadi $x = 0$ atau $x = 5$ (jangan
sekali-kali membagi dengan $x$: penyelesaian $0$ akan hilang).

$(2x+1)^2 = (x-4)^2$: pindahkan ke kiri lalu faktorkan selisih dua
kuadratnya:
\[
(2x+1)^2 - (x-4)^2
= \bigl[(2x+1) + (x-4)\bigr]\bigl[(2x+1) - (x-4)\bigr]
= (3x - 3)(x + 5).
\]
Aturan hasil kali nol: $x = 1$ atau $x = -5$.
\end{solution}

\begin{solution}{exo:g10:algebra:8}
$(x+2)(4-x) > 0$: faktor $x + 2$ lenyap di $-2$ (positif sesudahnya),
$4 - x$ lenyap di $4$ (positif sebelumnya). Hasil kalinya positif
tepat ketika kedua faktornya positif, yaitu pada $\intoo{-2}{4}$.

$\frac{2x-6}{x+1} \geq 0$: pembilang nol di $3$, penyebut nol di
$-1$ (terlarang). \omterm{def:g10:algebra:signtable}{Tabel tanda}: pecahan itu positif ketika kedua bagiannya
bertanda sama, yaitu untuk $x < -1$ atau $x > 3$; nilainya nol di $x = 3$.
Himpunan penyelesaiannya: $\intoo{-\infty}{-1} \cup \intco{3}{+\infty}$.
\end{solution}

\begin{solution}{exo:g10:algebra:9}
Setelah $x$ bulan, layanan pertama berbiaya $9x + 15$ dan yang kedua $12x$.
Yang pertama lebih murah apabila
\[
9x + 15 < 12x
\ \Longleftrightarrow\
15 < 3x
\ \Longleftrightarrow\
x > 5 .
\]
Mulai bulan ke-$6$, layanan pertama menjadi pilihan yang lebih murah.
\end{solution}

\begin{solution}{exo:g10:algebra:10}
Nilai terlarang: $x = 1$. Untuk $x \neq 1$, kalikan kedua ruas dengan $x - 1$:
\[
x + 3 = 2(x - 1) = 2x - 2
\ \Longleftrightarrow\
x = 5 .
\]
Karena $5 \neq 1$, penyelesaiannya adalah $x = 5$. Periksa:
$\frac{5+3}{5-1} = \frac84 = 2$.
\end{solution}

\begin{solution}{exo:g10:algebra:11}
$(n+1)^2 - n^2 = n^2 + 2n + 1 - n^2 = 2n + 1$, yang ganjil untuk setiap
\omterm{def:g10:numbers:sets}{bilangan bulat} $n$.

Dua bilangan ganjil berurutan dapat ditulis $2n - 1$ dan $2n + 1$. Maka
\[
(2n+1)^2 - (2n-1)^2
= \bigl[(2n+1)+(2n-1)\bigr]\bigl[(2n+1)-(2n-1)\bigr]
= 4n \times 2 = 8n,
\]
yaitu kelipatan $8$.
\end{solution}

\begin{solution}{pb:g10:algebra:1}
\textbf{1.} $x^2 - 10x + 25 = (x - 5)^2$;
$9x^2 - 49 = (3x - 7)(3x + 7)$;
$x^2 + 6x + 9 - 4 = (x + 3)^2 - 2^2 = (x + 1)(x + 5)$.

\textbf{2.} $x = -1$ atau $x = -5$; \quad $x = \frac73$ atau
$x = -\frac73$.

\textbf{3.} $(5 - t)(5 + t) = 25 - t^2 = 21$ memberi $t^2 = 4$,
$t = 2$: kedua bilangan itu $3$ dan $7$ (jumlah $10$, hasil kali
$21$: cocok).

\textbf{4.} $(7 - t)(7 + t) = 49 - t^2 = 33$: $t = 4$, bilangannya
$3$ dan $11$.

\textbf{5.} $\left(\frac s2 - t\right)\left(\frac s2 + t\right)
= \left(\frac s2\right)^2 - t^2 = p$ memaksa
$t^2 = \left(\frac s2\right)^2 - p$. Penyelesaian ada tepat
ketika nilai ini $\geq 0$, yaitu ketika
$p \leq \left(\frac s2\right)^2$ --- hasil kali dua bilangan
yang jumlahnya tertentu tidak pernah melampaui kuadrat setengah
jumlahnya (soal pagar Ratu Dido di jilid sebelumnya sudah tahu).

\textbf{6.} $(x - a)(x - b) = x^2 - (a + b)x + ab$: jadi
\omterm{def:g10:algebra:equation}{persamaan} $x^2 - sx + p = 0$ berkata persis ``$x$ adalah salah satu
dari dua bilangan yang berjumlah $s$ dan berhasil kali $p$''.

\textbf{7.} $(x-5)^2 - 4 = x^2 - 10x + 25 - 4 = x^2 - 10x +
21$. Selisih dua kuadrat:
$(x - 5 - 2)(x - 5 + 2) = (x - 7)(x - 3) = 0$: penyelesaiannya $3$
dan $7$ --- bilangan sang juru tulis pada pertanyaan 3.

\textbf{8.} $x^2 + 6x - 7 = (x + 3)^2 - 16 = 0$ memberi
$(x + 3)^2 = 16$, $x + 3 = \pm 4$: $x = 1$ atau $x = -7$.

\textbf{9.} $x^2 + 4x + 5 = (x + 2)^2 + 1 \geq 1 > 0$ untuk setiap
$x$: tidak ada penyelesaian real. Dalam bentuk $(x + h)^2 + k$:
penyelesaian ada tepat ketika $k \leq 0$ (lalu $(x+h)^2 = -k$ mempunyai
dua akar $-h \pm \sqrt{-k}$); ketika $k > 0$ bentuk itu tetap
positif.

\textbf{10.} Persegi bersisi $x$ (luas $x^2$) dengan dua
persegi panjang $3 \times x$ (seluruhnya $6x$) membentuk huruf L berluas
$x^2 + 6x = 7$; pojok $3 \times 3$ yang hilang (luas $9$)
melengkapkan sebuah persegi besar bersisi $x + 3$. Menambahkan $9$ pada
kedua ruas: $(x + 3)^2 = 16$, jadi $x + 3 = 4$ (panjang selalu positif)
dan $x = 1$. Aljabar dan gambarnya adalah satu tindakan yang sama.

\textbf{11.} $(x-3)(x+2)$: negatif di antara kedua akarnya. Penyelesaian
$< 0$: $x \in \intoo{-2}{3}$.

\textbf{12.} $x^2 - 10x + 21 \leq 0$, yaitu
$(x - 3)(x - 7) \leq 0$: di antara kedua akarnya, ujungnya termasuk:
$x \in \intcc{3}{7}$.

\textbf{13.} Nol di $x = 1$, terlarang di $x = -2$; pecahan
itu positif di luar kedua akarnya:
$x \in \intoo{-\infty}{-2} \cup \intco{1}{+\infty}$.

\textbf{14.} $25 - (x - 5)^2 = 25 - x^2 + 10x - 25 =
x(10 - x)$: identitasnya terbukti. Karena $(x-5)^2 \geq 0$ selalu,
$A(x) \leq 25$, dengan kesamaan tepat ketika $x = 5$: kandang persegi
$5 \times 5$, luas maksimum $25$ m$^2$ --- tabel di jilid
sebelumnya, kini menjadi bukti dua baris.

\textbf{15.} Untuk $x > 0$, kalikan pernyataannya dengan $x$:
$x^2 + 1 \geq 2x$, yaitu $x^2 - 2x + 1 = (x - 1)^2 \geq 0$ ---
benar, dengan kesamaan hanya di $x = 1$. Membaginya kembali dengan
$x > 0$ mempertahankan semuanya (\cref{prop:g10:algebra:ineqrules}).

\textbf{16.} Perkalian silang pada $\frac1x = \frac{x}{1 - x}$
(semua panjang positif): $x^2 = 1 - x$, yaitu
$x^2 + x - 1 = 0$.

\textbf{17.} $x^2 + x - 1 = \left(x + \frac12\right)^2 -
\frac54 = 0$ memberi $x + \frac12 = \pm\frac{\sqrt5}{2}$. Akar
positifnya: $x = \frac{\sqrt5 - 1}{2} \approx 0.618$.

\textbf{18.} $\varphi = \frac{2}{\sqrt5 - 1} =
\frac{2(\sqrt5 + 1)}{(\sqrt5 - 1)(\sqrt5 + 1)}
= \frac{2(\sqrt5 + 1)}{5 - 1} = \frac{\sqrt5 + 1}{2}
\approx 1.618$.

\textbf{19.} $\varphi^2 = \frac{(\sqrt5 + 1)^2}{4} =
\frac{6 + 2\sqrt5}{4} = \frac{3 + \sqrt5}{2}$, sedangkan
$\varphi + 1 = \frac{\sqrt5 + 3}{2}$: keduanya sama. Lalu
$\varphi - 1 = \frac{\sqrt5 - 1}{2} = x = \frac1\varphi$ (memang
begitulah $\varphi$ didefinisikan). Rasio emas adalah bilangan
yang kuadratnya bertambah satu dan kebalikannya berkurang satu.

\textbf{20.} $\sqrt5$ bersifat \omterm{ex:g10:numbers:classify}{irasional} (soal akhir pekan tentang sifat
\omterm{ex:g10:numbers:classify}{irasional} di jilid sebelumnya: $5$ bukan kuadrat sempurna), dan menambah
$1$ lalu membagi dua mempertahankan sifat \omterm{ex:g10:numbers:classify}{irasional} itu (operasi rasional,
\cref{pb:g10:numbers:1}).
Rasionya: $\frac53 \approx 1.667$, $\frac85 = 1.600$,
$\frac{13}{8} = 1.625$, $\frac{21}{13} \approx 1.615$:
berganti-ganti di atas dan di bawah, sambil mengerucut ke
$\varphi = 1.618\ldots$ --- bilangan pencacah irama pada soal akhir pekan
tentang pencacahan irama di jilid sebelumnya diam-diam mengejar potongan
emas; bab tentang barisan akan menangkapnya beraksi.

\end{solution}
